\begin{helper}
In the setting of the split-outside-blocker comparison, assume the source and
target activation--priority laws have the Bernoulli activation atom formulas
and the finite priority lower-rectangle product formulas with
\(0\le \gamma/\Delta\le 1\).  Let \(r\in V\), \(X\subseteq N_G(r)\), and let
\(\phi:V\to V'\) and \(\beta\) satisfy the local-copy and private-blocker
hypotheses used in the replacement construction: \(\phi\) is injective and
preserves all adjacencies inside \(X\cup\{r\}\); the vertices
\(\beta(x,z)\), indexed by outside incidences
\[
  x\in X,\qquad z\notin X\cup\{r\},\qquad xz\in E(G),
\]
are distinct, fresh from \(\phi(X\cup\{r\})\), adjacent only to the intended
\(\phi(x)\) among the embedded candidates, nonadjacent to \(\phi(r)\), and
exhaust all non-local neighbors of every \(\phi(x)\), \(x\in X\).
Then there is one common witness system for the actual hybrid comparison.  It
consists of a finite enumeration \(z_0,\ldots,z_{n-1}\) of the outside
blockers, candidate sets \(B_k=\{x\in X:xz_k\in E(G)\}\), source common and
extension coordinates, target common and extension coordinates, a single
finite atom index set \(\Omega\), paired source and target cells
\(C_\omega,C'_\omega\), a common nonnegative weight \(w(\omega)\), complete
source and target mixed-stage events \(S_j(\omega)\) and \(T_j(\omega)\), and
normalized masses \(P(j,\omega)\).  The coordinates satisfy the actual
enumeration and transport clauses: the enumeration is injective, the source
common coordinates contain
\(X\cup\{r\}\), the source extension coordinates contain all \(z_k\) and are
disjoint from the common coordinates, the target common coordinates are the
\(\phi\)-image of the source common coordinates, and the target extension
coordinates contain all private blockers \(\beta(x,z_k)\) with \(x\in B_k\).

The same cells \(C_\omega,C'_\omega\) simultaneously have the following
properties.  They are pairwise disjoint on each side, cover the corresponding
unit-priority support, and the complements of the unit-priority supports have
zero measure.  Their weights agree with the source and target cell masses:
\[
  \operatorname{ofReal}(w(\omega))=\nu(C_\omega)
  =\nu'(C'_\omega),
  \qquad w(\omega)\ge0 .
\]
For every \(j\le n\), the complete source stage \(S_j(\omega)\) is the event
that some \(x\in X\) survives the embedded local-neighborhood tests and all
retained common blockers \(z_i\) with \(j\le i<n\); the complete target stage
\(T_j(\omega)\) is the event that the corresponding \(\phi(x)\) survives the
embedded local-neighborhood tests and all already-replaced private blockers
\(\beta(x,z_i)\) with \(i<j\).  The registered masses \(P(j,\omega)\) are
nonnegative and vanish on zero-weight atoms.  They are the values of a
hybrid/composite stage chain used for endpoints and adjacent monotonicity,
not same-index identifications of the source and final-target projections.
The endpoint values represent the true source event at \(j=0\) and the true
target event at \(j=n\), and the family is monotone across adjacent
replacements:
\[
  P(k,\omega)\le P(k+1,\omega) \qquad(k<n).
\]
In addition, these endpoint identities aggregate over the same finite atom
family:
\[
  \operatorname{ofReal}\!\left(\sum_{\omega\in\Omega}w(\omega)P(0,\omega)\right)
  =
  \nu\{(A,\pi):I_G(A,\pi)\cap X\ne\emptyset\},
\]
\[
  \operatorname{ofReal}\!\left(\sum_{\omega\in\Omega}w(\omega)P(n,\omega)\right)
  =
  \nu'\{(A',\pi'):I_{G'}(A',\pi')\cap\phi(X)\ne\emptyset\}.
\]
If \(n=0\), these two displayed endpoint identities are both asserted
directly for the same value \(P(0,\omega)=P(n,\omega)\); no adjacent step is
available or used to relate them.
\end{helper}
\begin{proof}
Apply
\(\noderef{SamplingSplitOutsideBlockerActualSameSurfaceInputConstruction}\)
to the two activation--priority laws, the ratio hypotheses
\(0\le\gamma/\Delta\le1\), the four activation and priority product
formulas, and the replacement data \(\phi,r,X,\beta\).  The hypotheses passed
to it are exactly the local-copy and private-blocker hypotheses in the
assembly statement: \(X\subseteq N_G(r)\), injectivity of \(\phi\),
adjacency preservation on \(X\cup\{r\}\), injectivity and freshness of
\(\beta\), intended adjacency to \(\phi(x)\), nonadjacency to \(\phi(r)\), no
adjacency to another embedded candidate \(\phi(y)\) unless \(y=x\), and
exhaustion of all nonlocal neighbors of every \(\phi(x)\).  Its statement returns
one actual witness surface:
\[
  n,\ z_0,\ldots,z_{n-1},\ B_k,\ S,S_{\rm ext},\ S',S'_{\rm ext},T',
  \Omega,\ C_\omega,\ C'_\omega,\ w(\omega),
\]
together with a proof \(h_{\rm out}(k)\) that every listed \(z_k\) lies
outside \(X\cup\{r\}\).  The same node gives the injectivity of
\(k\mapsto z_k\), the assertion that each listed \(z_k\) is adjacent to some
candidate in \(X\), the converse coverage of every outside blocker adjacent
to \(X\), and
\[
  B_k=\{x\in X:xz_k\in E(G)\}.
\]
It also gives all source and target coordinate clauses asserted in this
assembly: \(X\cup\{r\}\subseteq S\), \(z_k\in S_{\rm ext}\), disjointness of
the source common and extension coordinates, \(S'=\phi(S)\), disjointness of
the target common and extension coordinates, \(T'=S'\cup S'_{\rm ext}\),
\(\phi(X\cup\{r\})\subseteq S'\), and membership of every relevant private
blocker \(\beta(x,z_k)\) in \(S'_{\rm ext}\).

The construction node also returns the prescribed common refined-cell
clauses on exactly the same surface.  Hence, for every \(\omega\),
\[
  w(\omega)\ge0,\qquad
  \operatorname{ofReal}(w(\omega))=\nu(C_\omega),\qquad
  \operatorname{ofReal}(w(\omega))=\nu'(C'_\omega).
\]
The source cells are pairwise disjoint, the target cells are pairwise
disjoint, the cells cover the two unit-priority supports, the complements of
the two unit-priority supports have zero measure, and membership in
\(C_\omega\) is transported to membership in \(C'_\omega\) by equality of the
activation trace and priority values on the common coordinates.  Finally, and
crucially, the same construction node supplies a
\(\noderef{SamplingSplitOutsideBlockerActualSameSurfaceInputData}\) package
on this same \(n,z,B,\Omega,C,C',w\) surface.  This package is the paired
mixed surface together with the source endpoint identity, the target endpoint
identity, and the adjacent shared-context comparisons.  No second blocker
enumeration, atom type, cell family, or weight function is introduced.

Define the source complete stage \(S_j(\omega)\) by the displayed formula in
the statement: a state lies in \(S_j(\omega)\) when some \(x\in X\) is active,
beats every active embedded neighbor from \(X\cup\{r\}\), and is not killed by
any retained common blocker \(z_i\) with \(j\le i<n\) and \(x\in B_i\).
Define the target complete stage \(T_j(\omega)\) by the displayed transported
formula: \(\phi(x)\) is active, beats active embedded neighbors, and is not
killed by any already installed private blocker \(\beta(x,z_i)\) with
\(i<j\) and \(x\in B_i\).  These definitions give the two stage-description
clauses of the assembly directly.  They also fix the indexing convention
used throughout the hybrid chain: source blockers with index at least the cut
remain common, while target blockers with index below the cut have been
split into private coordinates.

Use the same-witness adjacent-event node
\(\noderef{SamplingSplitOutsideBlockerSameWitnessAdjacentEvents}\) on these
very cells and blocker sets to name the current left and right boundary
events.  Its statement gives the displayed left boundary with all retained
source blockers except the current one, and the displayed right boundary
with all already-private target blockers except the current one, using a
flattened notation for the current private blocker.  Therefore the adjacent
events are still attached to the same \(n,z_k,B_k,\Omega,C_\omega,C'_\omega\)
surface.

Unpack the supplied same-surface input package.  Its surface field is a
paired mixed measure surface \(\mathcal S\), its endpoint fields identify
\(\mathcal S.M_0(\omega)\) and \(\mathcal S.M_n(\omega)\) with the true
source and target endpoint masses in \(C_\omega\) and \(C'_\omega\), and its
adjacent field states that consecutive mixed masses split as
\[
  \mathcal S.M_k(\omega)=U+L,\qquad
  \mathcal S.M_{k+1}(\omega)=U+R,\qquad L\le R .
\]

Now apply
\(\noderef{SamplingSplitOutsideBlockerActualCompleteStageMassLaw}\) to the
same graphs, measures, graph-replacement data, enumeration, coordinate
sets, cells, weights, just-defined stage events, adjacent boundary events,
and paired-input data \(\mathcal S\).  The hypotheses of that node are exactly
the coordinate clauses from
\(\noderef{SamplingSplitOutsideBlockerActualCoordinatePackage}\), the common
cell clauses from
\(\noderef{SamplingSplitOutsideBlockerPrescribedCommonRefinedCellPackage}\),
the two displayed definitions of \(S_j(\omega)\) and \(T_j(\omega)\), the
adjacent-event descriptions from
\(\noderef{SamplingSplitOutsideBlockerSameWitnessAdjacentEvents}\), and the
paired surface, endpoint, and adjacent shared-context premises supplied in
the preceding paragraph.  The mass law supplies a single family
\(P(j,\omega)\) of real normalized complete-stage values with
\[
  P(j,\omega)\ge0,\qquad w(\omega)=0\Rightarrow P(j,\omega)=0,
\]
the true source endpoint identity at \(j=0\), the true target endpoint
identity at \(j=n\), and the adjacent monotonicity
\[
  P(k,\omega)\le P(k+1,\omega)\qquad(k<n).
\]
Summing the endpoint cell identities over the same finite atom family gives
the two aggregate endpoint identities asserted in the statement, as follows.
For the source endpoint event
\[
  E=\{(A,\pi):I_G(A,\pi)\cap X\ne\emptyset\},
\]
take the support \(U=\bigcup_{\omega\in\Omega} C_\omega\).  The refined-cell
package supplies measurability of \(E\) and of every \(C_\omega\), and the
complete-stage mass law supplies
\[
  \operatorname{ofReal}(w(\omega)P(0,\omega))=\nu(E\cap C_\omega).
\]
The containment \(C_\omega\subseteq U\) is immediate from the definition of
\(U\), and
\[
  E\cap U=\bigcup_{\omega\in\Omega}(E\cap C_\omega)
\]
is the elementary distributivity of intersection over this finite union.  To
check the null-complement hypothesis, observe that if a source state has all
priority coordinates in \([0,1]\), the refined-cell package puts it in some
\(C_\omega\).  Hence \(E\setminus U\) is contained in the complement of the
unit-priority support, whose \(\nu\)-measure is zero by the same package.  The
other hypotheses of
\(\noderef{SamplingFiniteCellEndpointSumDecomposition}\) are exactly the
nonnegativity of \(w\), nonnegativity of \(P(0,\omega)\), the zero-weight
convention, and pairwise disjointness of the cells.  Applying that node gives
\[
  \operatorname{ofReal}\!\left(\sum_{\omega\in\Omega}w(\omega)P(0,\omega)\right)
  =\nu(E).
\]
The target endpoint is identical: take
\[
  E'=\{(A',\pi'):I_{G'}(A',\pi')\cap \phi(X)\ne\emptyset\},\qquad
  U'=\bigcup_{\omega\in\Omega} C'_\omega,
\]
use the target measurability, target cell disjointness, target unit-support
cover, target null-complement clause, and the target atomwise identity
\[
  \operatorname{ofReal}(w(\omega)P(n,\omega))=\nu'(E'\cap C'_\omega),
\]
then apply
\(\noderef{SamplingFiniteCellEndpointSumDecomposition}\) again.  This proves
the two finite weighted endpoint identities in the assembly statement.

Thus all asserted witnesses and all asserted properties live on one common
\(n,z_k,B_k,\Omega,C_\omega,C'_\omega,w\) surface.  The proof uses the
hybrid mass law as the actual supplier of the complete-stage family; it does
not infer the family from the algebraic hybrid producer, and it does not
assert any same-index source--target stage equality or any primitive
outside-\(B_k\) paired atom equality.
\end{proof}
